\subsubsection{Masa y resorte en el borde de una piscina}

Considere un sistema compuesto por un resorte y una masa que se encuentran al borde de una piscina muy profunda, como se indica en la Figura \ref{fig:piscina}. El resorte es de largo natural $\ell_0$ y constante elástica $k$. A éste se fija una pared móvil de masa despreciable. El sistema se prepara de tal modo que la partícula puntual de masa $m$ se coloca junto a esta pared en su posición de compresión máxima, es decir en $x = -\ell_0$, y se suelta desde el reposo.

\begin{enumerate}
    \item ¿Cuál es la condición que asegura que la masa $m$ se moverá desde $x = -\ell_0$?
    \item Encuentre el valor máximo de $\mu_d$ que permita a la masa llegar al borde de la piscina ($x = 0$) con velocidad no nula. Entregue el valor de esta velocidad.
    \item Una vez que entra a la piscina ($x > 0$), la masa experimenta una fuerza de roce viscoso lineal de constante $\gamma$. Suponga que no hay fuerza de empuje. Determine el alcance máximo y su velocidad límite.
\end{enumerate}

\begin{figure}[htbp]
    \centering
    \begin{tikzpicture}[>=Latex, scale=1.5]
        % --- Parámetros ---
        \def\Lzero{1.8} % Longitud L0
        
        % --- 1. El Suelo y la Piscina ---
        % Suelo (izquierda)
        \fill[gray!20] (-2.5, -2) rectangle (0, 0);
        \draw[thick] (-2.5, 0) -- (0, 0);
        \draw[pattern=north east lines, pattern color=gray!50] (-2.5, -2) rectangle (0, 0);
        
        % Piscina (derecha)
        \fill[blue!5] (0, -2) rectangle (2.5, 0);
        \draw[blue!40] (0, 0) -- (2.5, 0);
        % Líneas de agua
        \foreach \y in {-0.3, -0.7, -1.1, -1.5} {
            \draw[blue!30, thin] ({0.2-0.1*rand}, \y) -- ({2.3+0.1*rand}, \y);
        }
        \node[blue] at (0.5, -1) {$\gamma$};

        % --- 2. Ejes de Coordenadas ---
        \draw[->, thick] (-2.5, 0) -- (2.5, 0) node[right] {$x$};
        \draw[->, thick] (0, 0) -- (0, 1.2) node[above] {$y$};

        % --- 3. El Mecanismo (Muelle y Masa) ---
        % Pared fija
        \fill[gray!60] (-2.5, 0) rectangle (-2.2, 0.5);
        \draw[thick] (-2.2, 0) -- (-2.2, 0.5);
        
        % Resorte
        \draw[decoration={aspect=0.3, segment length=3pt, amplitude=3pt, coil}, decorate] 
            (-2.2, 0.25) -- (-\Lzero+0.2, 0.25);
        
        % Masa P
        \shadedraw[ball color=gray!80] (-\Lzero+0.2, 0.25) circle (7pt);
        \filldraw (-\Lzero+0.2, 0.25) circle (1pt); % centro
        
        % Cota L0
        \draw[|<->|] (-2.2, 0.7) -- (0, 0.7) node[midway, fill=white] {$\ell_0$};

        % --- 4. Traectoria de caída ---
        \draw[dashed, ->, gray] (0,0.15) parabola (1.8, -2);

        % --- 5. Etiquetas de coeficientes ---
        \node at (-1, -0.3) {$\mu_e, \mu_d$};
        \draw[->, thick] (1.2, 0.8) -- (1.2, 0.3) node[right] {$\mathbf{g}$};

    \end{tikzpicture}
    \caption{Sistema masa-resorte en el borde de una piscina con rozamiento seco y viscoso.}
    \label{fig:piscina}
\end{figure}